\section*{Chapitre \ref{ch:s2:the-bond-futures-basis-trade} --- Le basis trade sur obligations et contrats à terme}

\begin{solution}{exo:s2:the-bond-futures-basis-trade:1}
$1/(0.01 + 0.01) = 50$ fois~; à 3~\% chacun, $1/0.06 = 16.7$ fois.
\end{solution}

\begin{solution}{exo:s2:the-bond-futures-basis-trade:2}
$0.002 \times 50 = 10~\%$ par an sur le capital, avant frais et capitalisation.
\end{solution}

\begin{solution}{exo:s2:the-bond-futures-basis-trade:3}
$50 \times 0.006 = 30~\%$ du capital.
\end{solution}

\begin{solution}{exo:s2:the-bond-futures-basis-trade:4}
À dix fois, le capital couvre le collatéral sous stress (6~\% du nominal nécessitant 0.6 de capital), de sorte que le portefeuille conserve la totalité de sa position et récupère la base.
À cinquante fois, il doit vendre jusqu'à revenir à 16.7 fois sur la base stressée~; les positions vendues ne se redressent jamais, et les ventes coûtent 10 cents pour 100.
\end{solution}

\begin{solution}{exo:s2:the-bond-futures-basis-trade:5}
$(659 - 554)/659 = 15.9~\%$.
\end{solution}

\begin{solution}{exo:s2:the-bond-futures-basis-trade:6}
Les fonds à effet de levier détiennent des contrats à terme sur Treasuries pour de nombreuses raisons (vues sur la duration, trades de courbe, suivi de tendance), et le rapport ne montre que les contrats à terme,
et non les obligations au comptant ni le repo de l'autre côté~; une position vendeuse nette peut relever d'un basis trade ou d'un pari directionnel pur.
\end{solution}

\begin{solution}{exo:s2:the-bond-futures-basis-trade:7}
Le portefeuille à cinquante fois vend encore pendant 66 jours, puisqu'il se trouve au plafond de son collatéral~; son creux est plus profond ($-27.7~\%$) parce qu'il réduit moins tôt, et il
conserve davantage de la reprise ($-6.2~\%$ après, contre $-9.6~\%$). Une partie de la perte verrouillée vient des conditions de collatéral plus élevées,
l'autre du fait de fonctionner au plafond.
\end{solution}

\begin{solution}{exo:s2:the-bond-futures-basis-trade:8}
La convergence à la livraison ne profite qu'à un portefeuille encore détenu à la livraison. Les appels de marge et un repo qui ne peut pas être renouvelé peuvent forcer une vente avant
cette date, à la base du jour, et la perte devient alors réelle.
\end{solution}

\begin{solution}{pb:s2:the-bond-futures-basis-trade:1}
\begin{enumerate}
\item Base brute~: prix de l'obligation moins prix du contrat à terme multiplié par le facteur de concordance~; base nette~: base brute diminuée du carry~; taux de repo implicite~: le taux de financement auquel
l'opération d'achat, de livraison puis de remboursement s'équilibre.
\item Lorsque le taux de repo implicite dépasse le taux de repo payé.
\item Nominal par unité de capital, plafonné par l'inverse de la décote plus la marge~; \omterm{def:s2:the-bond-futures-basis-trade:rollover}{risque de renouvellement du repo}, c'est-à-dire le risque que le repo ne puisse pas être refinancé aux mêmes conditions.
\item Le collatéral~: décote plus marge sur contrats à terme.
\item Ils veulent de la duration sans mobiliser leur bilan et achètent des contrats à terme.
\item De 400 à 500 milliards de dollars au plus haut, plus de 60~\% de l'exposition des hedge funds sur Treasuries et 70~\% de leurs emprunts en repo.
\item Les fonds à effet de levier affichent une position vendeuse nette de 475 milliards de dollars au 18 février 2020, un record de 1\,184 milliards de dollars le 12 novembre 2024.
\item Ils détiennent des Treasuries en face de leurs positions acheteuses sur contrats à terme et les financent en repo.
\item Fourchettes plus larges, taux de repo qui bondissent, spreads d'arbitrage qui divergent, marges sur contrats à terme plus élevées.
\item Réduction des contrats à terme de 659 à 554 milliards de dollars, et vente pour 91 à 105 milliards de dollars de Treasuries.
\item Au moins 500 milliards de dollars d'achats de Treasuries le 15 mars, puis des achats aux montants nécessaires le 23 mars.
\item Non~: les Treasuries détenues dans l'opération ont continué de se négocier avec une prime avant le 17 mars.
\item \textbf{Résultat nommé.} Le portefeuille rapporte 2.1~\% par an à dix fois et 10.3~\% à cinquante fois~; dans le stress de financement simulé, le portefeuille à cinquante fois est contraint
de vendre, et termine 9.6~\% sous son capital d'avant stress une fois la base remontée~; le portefeuille à dix fois ne perd que 0.2~\%.
\item Valorisation $-6.6~\%$, coûts des ventes forcées $-3.3~\%$, repo plus cher $-0.6~\%$, carry $+0.8~\%$.
\item Le repo à terme maintient le financement tout au long du stress, de sorte que seule la marge sur contrats à terme peut forcer des ventes.
\item Positions sur contrats à terme des fonds à effet de levier, volumes de repo, écart entre le rendement implicite des contrats à terme et le rendement au comptant, conditions de collatéral.
\item De sorte que des conditions de collatéral sous stress et un mouvement de la base laissent le capital au-dessus de l'exigence.
\item Stresser les conditions de collatéral, inclure les coûts de vente forcée, utiliser le repo réellement disponible.
\item Le signal d'encombrement du basis trade.
\item Bilan~: il détient des Treasuries pour ceux qui veulent des contrats à terme.
\end{enumerate}
\end{solution}

\begin{solution}{iq:s2:the-bond-futures-basis-trade:1}
Les contrats à terme sont chers par rapport au comptant lorsque la demande à leur égard est forte~; un trader achète l'obligation moins chère à livrer (CTD), la finance en repo, vend le contrat à terme et perçoit
le taux de repo implicite moins le taux de repo, à fort effet de levier car les deux jambes sont collatéralisées.
\end{solution}

\begin{solution}{iq:s2:the-bond-futures-basis-trade:2}
Une hausse conjointe de la décote et de la marge, un mouvement du basis défavorable au portefeuille et une flambée du repo, appliqués à l'ensemble de la position d'un coup~; dimensionner le levier de sorte que
le capital couvre encore le collatéral dans les conditions stressées, et conserver un financement à terme.
\end{solution}

\begin{solution}{iq:s2:the-bond-futures-basis-trade:3}
À partir du prix plein de l'obligation aujourd'hui, du prix de facturation à la livraison (contrat à terme multiplié par le facteur de concordance, plus les intérêts courus) et des coupons intermédiaires~; un taux implicite
supérieur au taux de repo, les contrats à terme sont chers~; inférieur, ils sont bon marché.
\end{solution}

\begin{solution}{iq:s2:the-bond-futures-basis-trade:4}
Apporter davantage de capital si c'est possible~; sinon réduire la position, de préférence d'abord dans les jambes les moins liquides, et chercher un repo à terme~; ne jamais laisser
un appel de marge sans réponse.
\end{solution}

\begin{solution}{iq:s2:the-bond-futures-basis-trade:5}
Cours et taux de repo chaque matin~; exigences de collatéral de chaque prêteur et de la bourse~; marge de variation~; contrôle du capital face aux conditions
stressées~; calendriers de roll~; avec des alertes lorsque le capital s'approche de l'exigence.
\end{solution}

\begin{solution}{iq:s2:the-bond-futures-basis-trade:6}
Après un mouvement $x$, le capital vaut $C - Nx$ et l'exigence $Nk$, de sorte que les ventes commencent lorsque $C - Nx < Nk$. Avec $N = LC$, le plus grand mouvement sans vente est
$x < 1/L - k$~: pour $L = 50$ et $k = 0.02$, il est nul, ce qui explique pourquoi un portefeuille à son plafond vend les jours ordinaires.
\end{solution}
